\section{量子纠缠的数学描述}
\subsection{可分离态与纠缠态}
如果一个多粒子系统的态可以写成单粒子态的直积形式：
\[
\ket{\psi} = \ket{\psi_1} \otimes \ket{\psi_2}
\]
则称这个态为可分离态。否则，称为纠缠态。

\subsection{纠缠态的例子}
考虑两个量子比特的贝尔态：
\[
\ket{\Phi^+} = \frac{1}{\sqrt{2}}(\ket{00} + \ket{11})
\]
这个态不能分解为两个单量子比特态的直积，因此是纠缠态。

\subsection{纠缠的数学特性}
对于一般的两粒子态：
\[
\ket{\Psi} = r\ket{a_0b_0} + u\ket{a_0b_1} + s\ket{a_1b_0} + t\ket{a_1b_1}
\]

当 $rt \neq us$ 时，这个态就是纠缠态。纠缠态的一个重要特性是：对其中一个粒子的测量会立即影响另一个粒子的状态。

\subsection{量子计算的优势}
由于量子寄存器中的每个状态可以同时处于所有可能的状态，量子计算具有很强的并行处理能力。这种并行性使得量子算法在某些问题上比经典算法有指数级的加速。

\paragraph*{量子并行性}
量子并行性允许量子计算机同时处理多个计算路径：
\[
U(\alpha\ket{x} + \beta\ket{y}) = \alpha U\ket{x} + \beta U\ket{y}
\]

这种特性使得量子算法能够同时探索多个解空间，从而在某些问题上实现指数加速。